\chapter{Choosing tools for an accountable research assistant}
\label{exec14:l6}
\label{ch:la05-prototype}

\section{Define the research task before choosing the software}

The proposed causal assistant has three separable responsibilities. It helps
a person express a banking decision, invokes mathematical operations whose
applicability is explicit, and preserves the evidence needed to interpret the
answer. A model backend proposes text or structured objects. An orchestration
framework arranges calls and state transitions. A causal engine supplies a
particular graphical or numerical calculation. Their outputs have different
evidential meanings even when one software package exposes all three through
the same chat interface. The architecture should make those meanings visible
before selecting a vendor or adopting a framework.

Consider a lender comparing an offer with no offer for accounts entering
arrears. The declared population, twelve-month outcome, information available
at first eligibility, treatment contrast and cost convention determine the
question. A model may propose a graph, explain why a field is measured too
late or draft an analysis plan. A numerical service may compute a conditional
effect under the supplied graph and sampling law. The controller must prevent
the first activity from silently certifying the assumptions of the second.
The servicing comparison in Section~\ref{exec13:agent-architecture} develops
the substantive causal alternatives; here the question is which software
boundary preserves them.

Let an analysis state be
\begin{equation}
 s=(q,d,g,a,m,r,e),
 \label{exec14:l6-state}
\end{equation}
where $q$ is the decision contract, $d$ the data and observation definitions,
$g$ the graph or graph class, $a$ the assumptions and their statuses, $m$ the
method request, $r$ the result and $e$ its provenance. Each component has a
version identity. An external grant $u$ specifies permitted operations and
inputs. It is not a field that the model can authorize for itself. A proposed
transition $f$ should be admitted only if its preconditions hold for the
current state and grant. The output must retain the identities it actually
used. These are proposed application invariants, not properties supplied by
the word agent in a package description. \citep{exec14_l6_design}

This formulation also prevents an unhelpful choice between unrestricted
automation and manual work. Some transitions are mechanical: parsing a
versioned request, verifying source bytes or detecting a graph cycle. Others
need scientific judgment: whether an omitted variable is plausibly a common
cause, whether a policy's information set reflects actual availability, or
whether a transport assumption is defensible. An admissible workflow can
automate the first group and record uncertainty in the second. A fixed
controller can continue independent analyses after one method refuses a
request, provided it does not substitute a new target without an explicit
revision.

\section{Backend alternatives and what can actually be compared}

An offline fixture is the first useful backend. It returns a predetermined
proposal whose intended errors and ambiguities are known. It can establish
whether the controller preserves a PAG, refuses an unsupported policy and
binds a result to the right horizon. It cannot establish how often a language
model generates a correct proposal. The existing local prototype accepts
authored or recorded proposals and makes no live model calls. Recorded
provider/model labels are caller-supplied provenance, not independently
authenticated proof of generation. \citep{exec13_architecture_local,exec14_l6_local}

A prospective local model moves the inference computation to a selected
host. Qualification would bind weight files, tokenizer, inference software,
quantization and decoding settings, model terms and any remotely retrieved
material. A prospective hosted model binds a provider endpoint, model
revision or declared alias policy, request/response schema, data transfer
boundary, usage accounting and service contract. These are comparison
requirements. Neither local execution nor a commercial contract by itself
establishes causal competence, confidentiality of every component, or
permission to use a particular source. Current model rankings, prices and
retention guarantees are not inferred from generic framework examples.
\citep{exec14_l6_design}

\begin{longtable}{@{}p{.19\textwidth}p{.34\textwidth}p{.40\textwidth}@{}}
\caption{Backend choices for the proposed assistant. Except for the offline
controller evidence, these are prospective contracts, not measured services.}
\label{exec14:l6-backends}\\
\toprule Choice & Main benefit to investigate & Evidence needed before use\\\midrule
\endfirsthead
\toprule Choice & Main benefit to investigate & Evidence needed before use\\\midrule
\endhead
Authored or recorded proposal & Deterministic fault and semantic cases without
provider cost & Exact bytes, declared origin and an independent expected
disposition; no model-quality inference.\\
Local model & A controlled inference environment and inspectable deployment
inputs & Weight/tokenizer terms and hashes, runtime dependencies, outbound
access and logging, capacity, task performance and full operating cost.\\
Hosted model & A service interface with outsourced inference operations &
Exact service/model identity, allowed inputs, supported schema behavior,
usage and price evidence, retention and recovery contract.\\
Selective routing & Use additional model work only for eligible hard cases &
A frozen routing rule, measured conditional error and cost, comparable
information and an explicit unavailable route when no eligible backend exists.\\
Human or deterministic baseline & An interpretable reference for the value
of assistance & The same task definitions, tools and total repair opportunity;
record review time and failure dispositions.\\
\bottomrule
\end{longtable}

The common envelope should preserve the exact input and response bytes,
provider or local runtime identity, relevant settings, tool definitions,
retrieval identities, schema version and measured usage when available.
An adapter should distinguish malformed structure, an explicit backend
refusal, truncation, timeout and an unknown remote outcome. It should not
replace any of these with an empty successful causal result. A structurally
valid graph can still contain unsupported causal edges, so structural
validation and scientific adjudication remain separate.

Model substitution is therefore a reviewed change to the analysis system.
If a provider alias or local weight file changes, previously saved outputs
remain evidence of their original run. A new run should receive a new
execution identity, even if it preserves the same business question.
Conversely, keeping the same model while changing the target population
requires a new scientific identity. Requiring these two distinctions avoids
both needless recomputation of unrelated artifacts and unsafe reuse of
results for a changed question. The custody construction in
Section~\ref{exec14:artifact-custody} describes the narrower local byte and
context checks already observed.

\section{Three orchestration frameworks under inspected revisions}

LangGraph offers a state-graph and checkpoint interface. At the inspected
revision dated 7 October 2026, its project declares version 1.2.14 and Python
3.10 or later. The selected compile path accepts a checkpointer. Its interrupt
documentation is precise about resumption: the interrupted node starts again
and re-executes its preceding logic. That behavior is useful for a review
step, but an expensive or externally consequential action before the
interrupt may be repeated unless the application separately journals or
deduplicates it. The proposed reuse decision is a conditional Python
adapter for state orchestration, with scientific identities and side-effect
records owned by the application. Porting the entire framework is not
justified by the present proposal. The inspected root license is MIT;
dependency and deployment review remains separate. \citep{v12_la04_langgraph}

Pydantic AI offers typed tool interfaces and usage controls. The inspected
7 October 2026 revision declares Python 3.11 or later and a dynamic package
version; no exact release number is inferred from that field. Selected tool
fields support sequential execution and an approval condition, but their
inspected defaults do not require either. The application must configure
consequential tools explicitly. Request counts can be checked before a model
call, whereas token totals generally arrive with responses; an optional
pre-counting path has its own support and overhead. The selected cost check
warns when no cost was calculated and cannot enforce the stated monetary
limit in that case. The proposed decision is conditional reuse for typed
adapters, backed by an application-owned budget ledger and unknown-cost
policy. The MIT notice does not resolve every model or data term.
\citep{v12_la04_pydantic_ai}

AutoGen's inspected April 2026 snapshot declares AgentChat 0.7.5 and a pinned
core dependency. Its README describes maintenance mode with no new features;
that is a source declaration, not an independently verified evaluation of
any successor product. The selected AssistantAgent is stateful and expressly
not safe for concurrent calls on the same instance. It can execute returned
tool calls concurrently, and reflection can introduce another inference
call. Its save/load methods preserve model context; they do not record the
outcome of every external action. The code notice inspected for AgentChat
is MIT while the root documentation notice is CC-BY-4.0. The proposal retains
AutoGen as a comparative architecture and possible isolated adapter, but
does not make it the default controller or treat conversation persistence
as an exactly-once journal. \citep{v12_la04_autogen}

\begin{longtable}{@{}p{.20\textwidth}p{.35\textwidth}p{.38\textwidth}@{}}
\caption{Framework decisions concern selected source paths; none of these
frameworks has been executed or qualified by this comparison.}\\
\toprule Candidate & Useful inspected mechanism & Condition on adoption\\\midrule
\endfirsthead
\toprule Candidate & Useful inspected mechanism & Condition on adoption\\\midrule
\endhead
LangGraph & State graph, checkpoint and interrupt & Preserve immutable
scientific inputs; make replay of pre-interrupt work deliberate.\\
Pydantic AI & Tool schemas, control fields and usage checks & Configure
authorization/sequencing; reserve shared budgets; refuse unsupported unknown
cost when a hard ceiling is required.\\
AutoGen & Stateful model context and tool/reflection loop & Isolate mutable
instances; account for concurrency and reflection; review maintenance and
code/documentation terms.\\
Plain local controller & Fixed supported operations and a narrow journal &
Retain the observed limited scope; add capabilities only with their own
scientific and operating contracts.\\
\bottomrule
\end{longtable}

The plain controller is the baseline because its admitted operations and
failure states can be enumerated. This is not evidence that frameworks add
no value. A framework may reduce implementation effort for complex state,
streaming or human review. That advantage should be evaluated against the
additional dependencies, semantics and operational responsibilities it
introduces. A thin adapter can preserve the application's causal request
and result contracts while allowing this comparison. No framework should
be credited with discovery, identification or calibrated uncertainty merely
because it can call a tool bearing those names.

\section{Causal packages need decisions at the scientific boundary}

The package inventory covers different kinds of reusable material. PyWhyLLM
is a hypothesis-elicitation candidate: its inspected back-door suggestion
path delegates to a model, while the inspected front-door path is a
placeholder. A suggested confounder is not a proven adjustment set. Adopt
the idea of source-linked suggestions conditionally, and pass a supplied
graph to a qualified identification engine. Directly porting the entire
dependency stack would add maintenance without resolving that distinction.
\citep{v12_la_package_pywhy_llm}

The inspected CausalAgent package supplies interaction and tool-wrapper
ideas, but its OLC handler's missing function binding and permissive
malformed-output path prevent default direct adoption of that route. A
wrapper must carry typed errors and graph endpoint semantics. In oci-agent,
the useful intermediate objects are specifications, notebooks and diagnostic
reports. Its inspected unrestricted notebook runner has no finite per-cell
timeout, and its selected trimmed AIPW target differs from conventional
overlap weighting. Adapt the specification/report concepts behind bounded
operations; do not expose the general runner as a trusted causal endpoint.
These are static source decisions, not measured failure frequencies.
\citep{v12_la_package_causalagent,v12_la_package_oci_agent}

Benchmark software has another role. CLadder and CausalDS are candidates
for evaluating task formulations and answers after their generator/scorer
assumptions are reviewed. They are not engines whose test labels establish
the bank's causal truth. The selected IV-as-ATE and unknown-identification
boundaries remain reasons to quarantine unsupported cases or repair a
scorer before using it. OntoLearner can contribute structured vocabulary
and local versioned ontology objects, but taxonomy does not certify a
causal edge. LCM remains a deferred execution candidate until weights,
terms and its sufficient lagged regime are qualified. The LLM-CD and
LLM-SCD code pathways retain unresolved license or version-alignment
requirements; their papers can still be evaluated as scientific arguments.
\citep{v12_la_package_cladder,v12_la_package_causalds,v12_la_package_ontolearner,
v12_la_package_lcm,v12_la_package_llm_cd,v12_la_package_llm_scd}

These choices distinguish five useful dispositions. Direct adoption requires
the actual component's contract, environment and rights to match the use.
Conditional integration adds an adapter and predeclared acceptance cases.
Porting reimplements selected semantics and needs new numerical and error
comparisons; source similarity is insufficient. Reference-only use borrows
an idea without claiming executable reuse. Rejection or deferral records
the reason and a concrete reconsideration trigger. A package can receive
different dispositions for its graph representation, estimator and notebook
runner. Treating its repository as a single yes-or-no choice obscures the
component that creates the real scientific or operating risk.

\section{A state invariant for revision and invalidation}

Suppose the controller maintains a dependency graph of immutable analysis
objects. A result stores the exact identities of the target, input data,
graph, method and execution environment used to produce it. Define the
reuse invariant $I(s)$ to mean that every result presented as current has
matching admitted dependencies and retains its qualification status.
Initially, no result exists and $I$ holds. A transition that creates a
result establishes $I$ only after checking its dependencies. A transition
that changes a dependency preserves $I$ by retiring affected descendants
from the current view and requiring recomputation or an explicit stale
label. By induction, those transition rules preserve $I$ for a finite
sequence. This is an application-level argument under the stated rules,
not a proof that the current prototype implements an arbitrary dependency
graph or that supplied causal assumptions are true. \citep{exec14_l6_exact}

The retirement rule should be selective. In a finite illustrative dependency
graph, identification depends on the target and graph; estimation also
depends on data; valuation depends on the estimate and cost assumptions;
the report depends on valuation. A pure observation-clock diagnostic depends
on the data and clock rule. Changing the target horizon invalidates
identification, estimation, valuation and the report, while leaving that
diagnostic current if its own dependencies did not change. Changing costs
invalidates valuation and the report without changing the effect estimate.
An exact traversal verifies these sets in the project example. No claim is
made that every real diagnostic is target-independent: if its admissibility
rule uses the horizon, that dependency must be recorded too.
\citep{exec14_l6_exact}

This rule protects a banking decision from a subtle error. An agent can
produce a plausible revision to the question after observing a weak result.
The old estimate may remain numerically unchanged and its source files may
still pass all hashes. It becomes ineligible as an answer to the new
question because its scientific dependencies changed. The revised target
and the original result should both remain inspectable. The controller
must also prevent a later successful parse, retry or report step from
turning a conditional or unverified premise into a verified one. Such
promotion requires new supporting evidence and a separate recorded decision.

\section{Recovery includes an unknown outcome}

Local pure computations and external effects have different recovery
contracts. A deterministic calculation can be rerun with pinned inputs and
environment when recomputation is allowed. A paid model request may have
completed and been billed even if its response was lost. A proposed journal
distinguishes prepared, submitted, acknowledged, completed, failed and
outcome-unknown states. A timeout after submission should not be treated as
proof that the provider did no work. The availability of provider lookup or
idempotency must be established for the selected service.
\citep{exec14_l6_design}

There is a simple indistinguishability argument. Let history $h_0$ lose a
request before the external service receives it, and let history $h_1$
lose the response after the service completes the work. If the local
journal in each case contains only submitted and the same timeout, its
observable state is identical. A rule depending only on that state takes
the same action in both histories. Resubmission can duplicate work under
$h_1$; refusing to resubmit can leave the task unfinished under $h_0$.
Randomization changes the probabilities but cannot guarantee both no
duplication and eventual completion from this information alone. A service
lookup, idempotent external operation or an explicitly accepted duplication
policy adds information or changes the requirement. This finite argument
concerns the stated observation model, not all possible distributed
systems. \citep{exec14_l6_exact}

For a causal assistant, the appropriate response may be to retain an
unknown provider outcome while continuing independent deterministic checks
or another task. Cancellation similarly needs a scope: stopping future
local work does not prove that an already submitted remote request was
cancelled or its charge reversed. The budget ledger should retain any
reservation whose financial outcome is unknown until reconciliation. The
scientific report can remain incomplete with an explicit reason; filling
the missing answer from memory would sever its evidence chain.

\section{Shared budgets require reservations}

Usage checks on individual calls do not by themselves enforce a shared
budget under concurrency. Let $B$ be a stipulated expenditure ceiling,
$S$ settled cost and $R$ the sum of unresolved reservations. Before admitting
request $j$ with a credible upper cost bound $\bar c_j$, a controller can
require
\begin{equation}
 S+R+\bar c_j\le B.
 \label{exec14:l6-reserve}
\end{equation}
The check and reservation must be one serialized state transition. If the
actual cost $c_j$ is no larger than its reserved bound, completion replaces
that reservation with $c_j$; hence the total does not increase beyond its
admitted level. Starting from $S=R=0$, induction preserves $S+R\le B$.
Unknown prices, unbounded usage, excluded fees or provider overcharges
invalidate the premise; they are not resolved by this algebra.
\citep{exec14_l6_exact}

An exact illustrative case makes the race visible. Take $B=10$ dollars and
settled cost $S=6$. Two requests each require a three-dollar upper bound.
If each reads the same unreserved balance, both individually see
$6+3\le10$, yet their combined cost can reach 12 dollars. Under serialized
reservation, the first makes $S+R=9$, and the second is refused because
$6+3+3>10$. If the first settles at two dollars, $S$ becomes eight and $R$
returns to zero; the second three-dollar reservation is still inadmissible.
The project calculation enumerates both admission orders and verifies the
same conclusion. All amounts are stipulated; they are not prices of any
model. \citep{exec14_l6_exact}

A successful provider response may supply token counts after money has
already been spent. Accounting can record that cost correctly without
having prevented an overspend. Conversely, a token cap can limit one
resource while leaving human review time or a separately priced tool
unbounded. A useful prototype records call counts, known usage, reservations,
settled expense and unknown charges separately. Retries and reflection
consume the same shared budget, and failed requests remain in its total.
These proposed controls address the inspected framework boundaries without
claiming that a financial guarantee follows from a field named cost limit.

\section{Routing is a decision rule, not a model leaderboard}

A prospective router can compare eligible backends by expected task loss
and full incremental cost. For task information $z$, let $e_b(z)$ be a
calibrated error probability for backend $b$, $\ell(z)$ the loss of the
defined error, and $c_b(z)$ its expected call and review cost. Under the
deliberately simplified assumption of one relevant error category and a
fixed loss for that category, choose among eligible backends by
\begin{equation}
 b^*(z)\in\arg\min_{b\in\mathcal B(z)}
       \{c_b(z)+\ell(z)e_b(z)\}.
 \label{exec14:l6-route}
\end{equation}
The eligible set $\mathcal B(z)$ encodes data-use and capability constraints.
An ineligible backend cannot become eligible by having a low quoted cost.
If the set is empty, the workflow returns an unavailable route or uses an
authorized non-model alternative. Real error categories, uncertain losses
and abstention require a richer decision model. \citep{exec14_l6_exact}

For illustration, backend A costs one dollar with error probability .08,
and backend B costs four dollars with error probability .02. Then B has
lower expected total cost precisely when
$4+.02\ell<1+.08\ell$, or $\ell>50$ dollars. At loss 100, the stipulated
totals are nine and six dollars; at loss ten they are 1.80 and 4.20 dollars.
These numbers explain what a pilot should measure. They neither rank real
models nor justify applying a benchmark-wide error rate to a new banking
population. Costs of adjudication and wrong refusals must be included if
they differ between routes. \citep{exec14_l6_exact}

To evaluate a router, freeze its task strata and eligibility rules before
evaluation, give alternatives comparable information and tools, and retain
every assigned case. Selecting the backend after seeing the correct answer
would evaluate an oracle unavailable at deployment. A simple deterministic
form, a single assisted controller and a critic-assisted controller provide
useful alternatives before sophisticated routing. The experiments should
measure whether the extra model work changes target errors, unresolved
assumptions, usable results and review effort. Scientific identification
still belongs to the causal model and data; routing changes assistance,
not what the observational law identifies.

\section{From architectural comparison to a bounded prototype}

The proposed adoption sequence begins with the existing fixed-operation
controller and generated inputs. A selected framework can then be placed
behind an adapter preserving the same requests, grants, result types and
failure meanings. Offline fixtures test schema and recovery behavior before
an actual model is admitted. A live trial, if justified, needs named model
and service identities, permitted data, measured usage and an independent
task rubric. It can use public and generated tasks to establish whether
further bank experimentation is worth funding. No authentic customer
trial is required merely to compare the proposed interfaces.

The present comparison remains scoped. Three pinned orchestration
frameworks and the selected causal packages support concrete decisions;
they are not a complete census of all backends or hosted services. A
current commercial-model feature or price comparison requires inspected
official evidence and an applicable version/date. Model terms, deployment
rights, task performance, costs and authentic owner acceptance remain
distinct unresolved questions. The following preserved local implementation
provides a concrete baseline against which such additions can be tested.
\citep{exec14_l6_design,exec14_l6_local}

\label{exec14:l6-extension-end}
